\documentclass[11pt,a4paper]{article}
\usepackage[T1]{fontenc}
\usepackage[utf8]{inputenc}
\usepackage[polish]{babel}
\usepackage{amsmath,amssymb,amsthm,geometry,hyperref}
\geometry{margin=25mm}
\newtheorem{theorem}{Twierdzenie}
\newcommand{\Tr}{\operatorname{Tr}}
\newcommand{\Dtr}{D_{\rm tr}}
\newcommand{\CQ}{\mathcal C\!\mathcal Q}
\title{FIN: odporność dysonansu i identyfikowalność operacyjna}
\author{}
\date{}
\begin{document}
\maketitle

\begin{abstract}
Zakaz zerowego dysonansu kwantowego dla zadanego marginalnego strict
nie znika w pełnej rodzinie mikroskopowych oddziaływań zgodnych z tą
samą mieszaną dynamiką Hartreego. Pozorny wyjątek $a=-5/12$ zostaje
rozstrzygnięty dokładnym rachunkiem spektralnym. W większej klasie
zgodnej tylko z dynamiką stanów czystych otrzymujemy ilościowe
ograniczenie łączące nieklasyczność, asymetrię wymiany i brak stacjonarności.
Jednocześnie konstruujemy uniwersalne stacjonarne kanały o identycznych
obu lokalnych kanałach wyjściowych, lecz odmiennym splątaniu.
Ich mieszanka jest separowalna dokładnie przy równych wagach.
Wyniki określają, czego nie można wywnioskować z lokalnego jądra
i jakie dodatkowe pomiary lub zasoby rozróżniają opisy mikroskopowe.
Nie stanowią wyprowadzenia fundamentalnej teorii fizycznej FIN.
\end{abstract}

\section{Ustalone dane i trzy różne rodzaje informacji}

Na dwóch nośnikach o wymiarze $n=12$ przyjmujemy
\[
 V=\frac{|\Omega\rangle\langle\Omega|+S-2D}{2},\quad
 |\Omega\rangle=\sum_i|ii\rangle,\quad
 D=\sum_i|ii\rangle\langle ii|,\quad P_s=(I+S)/2,
\]
gdzie $S$ zamienia nośniki. Mapa źródłowa to
$T(\rho)=\Pi\Re\rho=\Tr_2[V(I\otimes\rho)]$.
\footnote{Źródła repozytoryjne:
\texttt{FIN\_Hartree\_Equivalence\_and\_Strict\_Correlation\_Floor.tex}
oraz \texttt{FIN\_Separable\_Stationarity\_and\_Necessary\_Discord.tex}.
Rozróżniają pełny operator źródłowy, przepływ stanów mieszanych
i słabszą zgodność czystych przepływów projektowych.}

Zadany rzeczywisty stan strict to $C=I/12+W/20$,
\[
 W_{ij}=\frac{\cos((743d_{ij}+650)/4000)}{1+d_{ij}^{9/5}},\qquad W_{ii}=0,
\]
z odległością cykliczną. Dla przepływu $i[T(\rho),\rho]$ wszystkie
hermitowskie podniesienia symetryczne względem zamiany, zgodne na
każdym stanie mieszanym, mają postać
\begin{equation}\label{eq:mixed}
 H_a=V+aS+bI.
\end{equation}
Zgodność tylko na czystych projektorach pozostawia większą klasę
\begin{equation}\label{eq:pure}
 H=V+cP_s+B_a,\qquad B_a=P_aB_aP_a=B_a^*,\quad P_a=(I-S)/2.
\end{equation}
Są to wcześniejsze klasyfikacje, nie nowe założenie, że obie przesłanki
są równoważne. Znajomość pełnej mapy $T$ jest jeszcze silniejsza.

Poniższe wyniki rozróżniają: lokalny przepływ, lokalne kanały
przygotowania oraz pełne wspólne rekordy pomiarowe. Nie są to te same
dane. Zamiana dwóch nośników nie jest kierunkiem generatora $\mathbb Z_{12}$,
więc żaden wniosek o asymetrii nie rozstrzyga selektora QW-2191.

\section{Dysonans jest konieczny w całej klasie mieszanego przepływu}

Cztery pasma $H_a-bI$ mają wartości
\[
 e_\Omega=a+(n-1)/2,\quad e_+=a+1/2,\quad
 e_D=a-1/2,\quad e_A=-a-1/2.
\]
Wartości kanonicznego $V$ na tych pasmach to odpowiednio
$(n-1)/2,1/2,-1/2,-1/2$.
Dla $j\in\{+,\Omega\}$ zdefiniujmy wielomiany Lagrange'a
$p_j(x)=\prod_{k\ne j}(x-e_k)/(e_j-e_k)$.
Wtedy
\begin{equation}\label{eq:poly}
 V=-I/2+p_+(H_a-bI)+(n/2)p_\Omega(H_a-bI)
\end{equation}
wszędzie poza $a=-1/2,-n/4$.
Zbieżność dwóch pasm przy $a=0$ nie szkodzi, bo mają tę samą wartość $V$.

Dla $n=12$ mianownik po uproszczeniu wynosi
$16(a+3)(2a+1)$. Jest niezerowy przy $a=-5/12$.
Dlatego każda równowaga w tym dawnym przypadku nierozstrzygniętym
jest także równowagą kanonicznego $V$.

Dwa rzeczywiste wyjątki wymagają osobnego dowodu. Istnieją tam dodatnie
stany stacjonarne z koherencją między pasmami, które nie komutują z $V$;
nie wolno więc rozszerzać \eqref{eq:poly} przez domysł.
Jednak blok między jednorodnym $u$ a naprzemiennym $v$ ma wartości
singularne $(n-2)/(2n)$, $|a+(n-2)/(2n)|$ i $1/n$.
Przy obu wyjątkach dla $n=12$ norma odwrotności wynosi 12, a normy
hamiltonianu, po pominięciu skalaru, wynoszą odpowiednio $5$ i $7/2$.
Nie przekraczają kanonicznego $11/2$.

\begin{theorem}[Jednolita granica dla wszystkich $a,b$]\label{thm:mixed}
Każdy stacjonarny stan dwuciałowy dla \eqref{eq:mixed}, którego
pierwszym marginalnym stanem jest zadany $C$, ma odległość od
zbioru klasyczno-kwantowego po tej stronie co najmniej
\begin{equation}\label{eq:d0}
 d_0=\frac{911271283377980905464649}
 {26400000000000000000000000000000}>3.45\cdot10^{-8}.
\end{equation}
\end{theorem}

\begin{proof}
Poza dwoma wyjątkami \eqref{eq:poly} przenosi stacjonarność na $V$.
W wyjątkach poprzedni dowód oparty na odwracalnym bloku zachowuje
tę samą normę odwrotności i nie większą normę hamiltonianu.
Niech $\delta_L$ będzie dokładną dolną granicą najmniejszej niezerowej
szczeliny widma $C$, a $\Delta_L$ dolną granicą różnicy jego prostych
wag modów $u,v$. Ten sam rachunek daje
$d_0=\delta_L\Delta_L/7392$.
\footnote{Przeliczane przedziały pochodzą z
\texttt{fin\_projected\_learning/research.py} i
\texttt{fin\_replication\_consistency/certify.py}.
W szczególności $\delta_L=27234855667/12500000000000$ oraz
$\Delta_L=234218204114629/2000000000000000$.}
Skalar $bI$ nie zmienia stacjonarności. Istnienie równowagi o tych
marginałach dla każdego $a$ nie jest częścią twierdzenia.
\end{proof}

Wniosek dotyczy zerowego jednostronnego dysonansu i geometrycznej
odległości od zbioru CQ. Nie podaje wartości entropowego dysonansu
ani miary splątania.

\section{Nieklasyczność, asymetria lub ruch: większa klasa czysta}

Dla dowolnego $R$ połóżmy $\overline R=(R+SRS)/2$.
W całej klasie \eqref{eq:pure} zachodzi dokładna tożsamość
\begin{equation}\label{eq:twirlidentity}
 [V,\overline R]=P_s[H,R]P_s.
\end{equation}
Po uśrednieniu stan jest blokowy względem wymiany; na bloku
symetrycznym $H$ różni się od $V$ tylko skalarem, a na
antysymetrycznym sam $V$ jest skalarny. To dowodzi tożsamości
także bez stacjonarności.

Niech $d_A(R)=\inf_{Q\in\CQ_A}\Dtr(R,Q)$,
$\mathcal A(R)=\Dtr(R,SRS)$ i $\epsilon=\|[H,R]\|_1$.
Zakładamy tutaj, że średnia obu oznakowanych marginałów jest równa $C$,
niekoniecznie każdy marginalny stan osobno.
\begin{theorem}[Trójczłonowe ograniczenie]\label{thm:tradeoff}
W powyższym zakresie
\begin{equation}\label{eq:tradeoff}
 d_A(R)+\tfrac12\mathcal A(R)
       +\frac{\delta_L}{616}\epsilon\ge d_0.
\end{equation}
Analogiczny wynik zachodzi po drugiej stronie.
\end{theorem}

\begin{proof}
Stan $\overline R$ ma oba marginały $C$. Rzutowanie go w różne
przestrzenie własne pierwszego $C$ daje $\widetilde R$.
Rachunek odwracalnego bloku i \eqref{eq:twirlidentity} dają
$c_0-c_6\le12\epsilon+24\|V\|\|\overline R-\widetilde R\|_1$.
Porównanie normy Frobeniusa i śladowej oraz dodatnia szczelina widma
prowadzą do
\[
 \|\overline R-\widetilde R\|_1
 \le\frac{48(1+\|C\|)}{\delta_L}\,d_A(\overline R).
\]
Ponieważ $\|V\|=11/2$, $\|C\|<1/6$, współczynnik łączny nie
przekracza $7392/\delta_L$. Odległość od dowolnego ustalonego zbioru
jest 1-lipschitzowska, więc
$d_A(\overline R)\le d_A(R)+\mathcal A(R)/2$.
Podzielenie nierówności przez ten współczynnik daje \eqref{eq:tradeoff}.
\end{proof}

Dla stanu stacjonarnego i symetrycznego dostajemy $d_A(R)\ge d_0$
niezależnie od 4357 niewidocznych parametrów oddziaływania.
Dla stacjonarnego stanu CQ konieczna byłaby natomiast
$\mathcal A(R)\ge2d_0$. Nie jest to dowód nasycenia tej granicy.
Dla hamiltonianu pomnożonego przez $g\ne0$ resztę komutatora należy
podzielić przez $|g|$; nie wyprowadzono absolutnej jednostki czasu.

Równe marginały same nie wymuszają symetrii wymiany. Nie wolno również
zakładać, że uśrednienie zachowuje CQ: dla małego $t>0$ stan
\[
 C\otimes C+t(P_u-P_v)\otimes(|u\rangle\langle v|+|v\rangle\langle u|)
\]
jest dodatni i CQ po pierwszej stronie, ma oba marginały $C$, ale
jego średnia po zamianie już nie komutuje z $C\otimes I$.
Ten przykład nie jest deklarowany jako stacjonarny.

\section{Jawny stan asymetryczny i dostępne rekordy}

Niech $C'=I/12+11W/100$; dokładnie $C'\succ I/125$.
Dla $Q_i=I-|i\rangle\langle i|$ zdefiniujmy
\begin{equation}\label{eq:cq}
 R_{\rm CQ}=\frac1{n-1}\sum_i |i\rangle\langle i|\otimes Q_iC'Q_i.
\end{equation}
Jest to wyjście kompletnego instrumentu Lüdersa
$K_i=Q_i/\sqrt{n-1}$, z zapisem wyniku na pierwszym nośniku.
Spełnia $\sum_iK_i^*K_i=I$ i jest deterministycznym kanałem CQ
na dowolnym wejściu. Alternatywna realizacja przez losowy wybór
projektora i selekcję sukcesu ma prawdopodobieństwo $11/12$;
nie należy utożsamiać tych dwóch realizacji sprzętowych.

Dla podanego programu marginały \eqref{eq:cq} wynoszą
\[
 (R_{\rm CQ})_A=I/12,\qquad
 (R_{\rm CQ})_B=I/12+W/10,
\]
a ich średnia jest $C$. To nie jest kontrprzykład z dwoma
indywidualnymi marginałami $C$.
Ponieważ $DR_{\rm CQ}=0$, hamiltonian
$H_-=V-S/2=(|\Omega\rangle\langle\Omega|-2D)/2$
anihiluje jego nośnik. Stan oraz każda z jego gałęzi są stacjonarne.
Średnia $\overline R$ jest natomiast równowagą kanonicznego $V$
z oboma marginałami $C$ i niezerowym dysonansem po obu stronach.
Stan \eqref{eq:cq} jest klasyczny tylko po jednej określonej stronie;
nie ustanawia w pełni klasycznej teorii dwóch nośników.

Dla efektów komutujących z $S$ mamy $\Tr(OR)=\Tr(O\overline R)$.
Co więcej, dla niekontrolowanych kanałów ewolucji
\begin{equation}\label{eq:passive}
 \mathcal P_S\mathcal U_a(t)=\mathcal U_0(t)\mathcal P_S,
 \qquad \mathcal P_S(R)=(R+SRS)/2.
\end{equation}
Każda wspólna, ustalona mapa wyniku instrumentu, kowariantna względem
$S$, komutuje z $\mathcal P_S$. Indukcja daje zatem identyczne rekordy
każdego skończonego adaptacyjnego protokołu zbudowanego z takich map
i pasywnych kanałów dla par opisów $(H_-,R_{\rm CQ})$ i $(V,\overline R)$.

To ograniczona równoważność operacyjna, nie zakaz rozróżnienia modeli.
Oznakowany lokalny pomiar widzi różne marginały. Sterowanie koherentne
hamiltonianem jest także mocniejszym zasobem niż jego zwykły kanał.
\footnote{M.~Araújo, A.~Feix, F.~Costa, Č.~Brukner,
\emph{Quantum circuits cannot control unknown operations},
New Journal of Physics 16, 093026 (2014),
\href{https://arxiv.org/abs/1309.7976}{arXiv:1309.7976}.
Ogólne rozróżnienie dostępu do bramki i jej kontrolowanej wersji
nie oznacza niemożliwości kontroli w określonej realizacji fizycznej.}
Z kontrolą początkowo w $|+\rangle$ odczyt $X$ wynosi
$\Re\Tr(U(t)R)$: tutaj jest równy $\cos(t/2)$ w pierwszym kanonicznym
opisie oraz $1$ w opisie przesuniętym. Przy $t=\pi$ otrzymujemy
$0$ i $1$. Oba kontrolowane hamiltoniany mogą respektować zamianę.
Dlatego sama symetria nie uzasadnia wykluczenia takiego dostępu.
Zapytania o energię i kontrolowane czasy nie należą do \eqref{eq:passive}.

\section{Te same lokalne kanały, różne splątanie}

Połóżmy $P_+=(I+S)/2-D$, $P_-=(I-S)/2$.
Dla dowolnego wejściowego stanu $\rho$ definiujemy
\begin{equation}\label{eq:channels}
 \mathcal E_\pm(\rho)=\frac2{n-1}P_\pm(\rho\otimes I)P_\pm.
\end{equation}
Kanały tego typu należy odróżnić od dokładnego klonowania.
Konstrukcje z kompresją do podprzestrzeni są znanym narzędziem
teorii przybliżonego klonowania; nie zgłasza się tu ich ogólnej optymalności.
\footnote{R.~F.~Werner, \emph{Optimal Cloning of Pure States},
Physical Review A 58, 1827 (1998), §3,
\href{https://arxiv.org/abs/quant-ph/9804001}{arXiv:quant-ph/9804001}.}

\begin{theorem}[Stacjonarna rodzina kanałów]\label{thm:channels}
Oba kanały \eqref{eq:channels} są CPTP, ich wyjścia są stacjonarne
dla $V$, a oba lokalne kanały są identyczne:
\begin{equation}\label{eq:local}
 \mathcal B(\rho)=\frac{I+(n-2)\rho}{2(n-1)}.
\end{equation}
Dla $\mathcal E_r=r\mathcal E_++(1-r)\mathcal E_-$,
$0\le r\le1$, wyjście jest separowalne dla każdego wejścia dokładnie
wtedy, gdy $r=1/2$. Przy każdym innym $r$ jest NPT dla każdego stanu wejściowego.
\end{theorem}

\begin{proof}
Operatory Krausa to
$\sqrt{2/(n-1)}P_\pm(I\otimes|j\rangle)$.
Ich zupełność wynika z $\Tr_2P_\pm=(n-1)I/2$.
Wyjścia należą do pasm energii $\pm1/2$ kanonicznego $V$, więc są
stacjonarne. Bezpośredni ślad częściowy daje \eqref{eq:local}.

Mieszanka z wagą $1/2$ jest dokładnie średnią po zamianie wyjścia
instrumentu \eqref{eq:cq}, z dowolnym $\rho$ zamiast $C'$.
Jest więc jawną sumą dodatnich produktów. Dla niezrównoważonej
mieszanki blok transpozycji częściowej na $|ii\rangle,|jj\rangle$
ma zerową przekątną i element pozadiagonalny
\[
 z=\frac{(2r-1)(\rho_{ii}+\rho_{jj})}{2(n-1)}.
\]
Istnieje para z $\rho_{ii}+\rho_{jj}>0$. Jej wyznacznik wynosi
$-|z|^2<0$, co dowodzi NPT i splątania.
\end{proof}

Dla jednolitej przekątnej wejścia negatywność jest co najmniej
$|2r-1|/n$, a $\Tr(S\mathcal E_r(\rho))=2r-1$.
Wspólny pomiar $S$ rozróżnia zatem tę rodzinę. Tylko lokalne kanały,
nawet badane z rejestrem odniesienia, są identyczne; nie wolno
przenosić tu pełnej równoważności pasywnej z poprzedniej sekcji.

Przy wejściu $C'$ oba lokalne wyjścia są dokładnie $C$.
Splątanie można więc zmienić bez zmiany całej lokalnej mapy przygotowania
i bez utraty stacjonarności. Przy wejściu $I/n$ te same kanały mogą
dawać splątanie mimo zerowego lokalnego źródła $T(I/n)$.
Splątanie samo nie określa ani nie generuje kształtu jądra.

Kanał $\mathcal E_+$ jest stacjonarny także dla całej klasy
\eqref{eq:pure}, ponieważ jego wyjście leży w ustalonym symetrycznym
paśmie. Zapewnia to niepusty przykład stacjonarnego uzupełnienia strict
niezależny od dowolnego bloku antysymetrycznego. Nie jest to jednak
wyprowadzenie programu $C'$ ani propagatora strict z samej stacjonarności.

\section{Koszt zasobów i granice interpretacji kanału}

Mieszanka z wagą $1/2$ jest deterministycznym globalnym kanałem CPTP
z jednego wejścia. Nie wolno na tej podstawie nazwać jej uniwersalnym
protokołem LOCC, gdy jedyny nieznany program znajduje się początkowo
u jednej ustalonej strony. Dla $n=12$ lokalny współczynnik depolaryzacji
wynosi $5/11$. Znormalizowany stan Choia ma po transpozycji częściowej
wartość własną
\[
 \frac{1-(n+1)(5/11)}{n^2}=-\frac3{88}.
\]
Kanał lokalny nie niszczy więc każdego splątania z odniesieniem.
Przesłanie tej odpowiedzi do drugiej strony wymaga transferu kwantowego
lub odpowiedniego wcześniejszego zasobu; losowy SWAP nie jest
automatycznie darmową operacją klasyczną.

Ograniczenie dotyczy uniwersalnego kanału na nieznanym wejściu,
potencjalnie splątanym z odniesieniem. Jeden ustalony znany stan
separowalny może być przygotowany lokalnie z jego klasycznego opisu.
Wcześniejsza konstrukcja LOCC z dwoma kopiami programu pozostaje
poprawną, inną architekturą zasobów.

Tożsamości kanałów nie zależą od konkretnego jądra. Dla osobnej
kanonicznej referencji legacy program
$\rho'=I/12+(11/5000)K_{\rm legacy}$ daje oba marginały
$I/12+K_{\rm legacy}/1000$. Dodatniość programu zapewnia wcześniejsze
oszacowanie $|K_{{\rm legacy},ij}|<3$, dające dolną granicę
$1/12-33(11/5000)=161/15000>0$.
Przenosi się więc tożsamość kanałów i klasyfikacja ich splątania,
nie liczba \eqref{eq:d0}, most między jądrami ani role fizyczne.
Wagi legacy o zmiennym znaku nie są tutaj klasycznymi szybkościami przejść.

Podstawowa teoria odniesień i ograniczeń operacyjnych jest znana;
zamiana dwóch abstrakcyjnych nośników nie jest tu ogłoszona
fizyczną regułą superselekcji ani źródłem kierunku w $\mathbb Z_{12}$.
Źródło musi ustalić, jakie obserwable, rekordy, połączenia kwantowe
i kontrolę czasu rzeczywiście ma dostępny obserwator.

\section{Bilans i dalszy test rozstrzygający}

\paragraph{Udowodniono.}
Wyjątek $-5/12$ nie znosi zakazu zerowego dysonansu; jednolita granica
obowiązuje w całej klasie tego samego mieszanego przepływu.
W większej klasie czystej zachodzi \eqref{eq:tradeoff}.
Istnieją dokładne stacjonarne kanały z tymi samymi lokalnymi odpowiedziami
i całkowicie sklasyfikowaną, odmienną zawartością splątania.

\paragraph{Obalono nadmierne uogólnienia.}
Symetria hamiltonianu i równe marginały nie wymuszają symetrii stanu.
Uśrednienie nie musi zachować CQ. Rekonstrukcji wielomianowej nie
wolno stosować w dwóch prawdziwych przecięciach pasm, choć osobne
kryterium nadal rozstrzyga nieklasyczność strict. Równoważność kanałów
pasywnych nie obejmuje sterowania hamiltonianem. Separowalne wyjście
nie czyni uniwersalnej mapy rozsyłania operacją LOCC bez zasobu kwantowego.

\paragraph{Pozostaje warunkowe.}
Konstrukcje korzystają z programu zawierającego $W$ i z jawnie zadanych
instrumentów. Klasyfikacja nie wyprowadza ich z nadsolitonu.
Przykład asymetryczny odtwarza tylko średnią marginałów, nie dwa
indywidualne strict. Istnienie stacjonarnego CQ z oboma marginałami
$C$ w całej klasie czystej bez symetrii stanu nie jest tu rozstrzygnięte.
Nie określono również optymalnych kosztów wszystkich protokołów.

Najbardziej rozstrzygające dalsze zadanie to źródło wspólnego kanału
przygotowania i dostępnej algebry operacyjnej. Lokalna geometria i jej
przepływ nie ustalają wagi $r$, wspólnych rekordów ani dostępu do
kontrolowanego hamiltonianu. Wspólny pomiar $S$ i kontrola fazy są
konkretnymi teoretycznymi sposobami rozróżnienia części modeli,
ale ich fizyczne pochodzenie i kalibracja pozostają otwarte.

Żaden wynik nie jest pomiarem laboratoryjnym, domknięciem QW-2191,
mostu legacy--strict, transferem ról, źródłem jednostek, SM/GR,
$L_{\rm total}$ ani ToE. Program i warunki początkowe mogą być
legalnymi wejściami fizyki, lecz nie wolno ich relabelować jako
wyprowadzonego samoistnie jądra. Nadsoliton pozostaje pierwotną
informacją w stanie solitonicznym, bez osobnej warstwy pod nim.

\section*{Źródła i dane}
\begin{enumerate}
\item Poprzednie raporty źródłowe FIN wskazane w pierwszym przypisie,
oraz dokładne przedziały w
\texttt{fin\_replication\_consistency/certify.py} i
\texttt{fin\_projected\_learning/research.py}.
\item M.~Araújo, A.~Feix, F.~Costa, Č.~Brukner,
\emph{Quantum circuits cannot control unknown operations},
New Journal of Physics 16, 093026 (2014),
\href{https://arxiv.org/abs/1309.7976}{arXiv:1309.7976}.
\item R.~F.~Werner, \emph{Optimal Cloning of Pure States},
Physical Review A 58, 1827 (1998),
\href{https://arxiv.org/abs/quant-ph/9804001}{arXiv:quant-ph/9804001}.
\item \texttt{fin\_discord\_robustness/research.py},
\texttt{test\_research.py}, \texttt{results.json} i
\texttt{verification.json}. Tożsamości symboliczne i wymierne
certyfikaty są odróżnione od pomocniczych kontroli numerycznych.
Znane narzędzia klonowania i kontroli kwantowej nie są zgłaszane
jako nowe ogólne prawa fizyki.
\end{enumerate}
\end{document}
